\documentclass[12pt]{article}
\usepackage[T1]{fontenc}
\usepackage[utf8]{inputenc}
\usepackage{lmodern}
\usepackage{textcomp}
\usepackage{enumitem}
\usepackage{geometry}
\usepackage{graphicx}
\usepackage{amsmath, amssymb}
\geometry{margin=1in}
\begin{document}
\begin{center}
Chemistry - Graduate Level Assessment 21 \
Total Marks: 40 \
Answer all questions.
\end{center}

\section*{Multiple Choice (10 marks)}
\begin{enumerate}[label=\arabic*.]
\item The executioner in charge of the lethal gas chamber at the state penitentiary adds excess dilute H\_2 SO\_4 to 196 g (about (1/2) lb) ofNaCN. What volume of HCN gas is formed at STP?
\begin{enumerate}[label=(\alph*)]
    \item 11.2 liters
    \item 67.2 liters
    \item 168 liters
    \item 100.8 liters
    \item 22.4 liters
\end{enumerate}
\item Calculate the average separation of the electron and the nucleus in the ground state of the hydrogen atom.
\begin{enumerate}[label=(\alph*)]
    \item 0.159 pm
    \item 79.5 pm
    \item 210 pm
    \item 132 pm
    \item 0.53 pm
\end{enumerate}
\item Calculate the number of H^+ ions in the inner compartment of a single rat liver mitochondrion, assuming it is a cylinder 1.8\ensuremath{\mu} long and 1.0\ensuremath{\mu} in diameter at pH 6.0.
\begin{enumerate}[label=(\alph*)]
    \item 1000 H^+ ions
    \item 900 H^+ ions
    \item 843 H^+ ions
    \item 950 H^+ ions
    \item 615 H^+ ions
\end{enumerate}
\item Which is correct about the calcium atom?
\begin{enumerate}[label=(\alph*)]
    \item It contains 20 protons and neutrons
    \item All atoms of calcium have 20 protons and 40 neutrons
    \item It contains 20 protons, neutrons, and electrons
    \item All calcium atoms have 20 protons, neutrons, and electrons
    \item All atoms of calcium have a mass of 40.078 u
\end{enumerate}
\item In which of the following is not the negative end of the bond written last?
\begin{enumerate}[label=(\alph*)]
    \item H-O
    \item P-Cl
    \item N-H
    \item S-O
    \item C-O
\end{enumerate}
\end{enumerate}

\section*{True / False (10 marks)}
\textit{Answer True or False}
\begin{enumerate}[label=\arabic*.]
\setcounter{enumi}{5}
\item True or False: Silver (Ag) possesses a higher electronegativity than Iron (Fe) due to its position in the periodic table.
\item True or False: The volume of the rectangular tank is 250 liters when calculated from its dimensions of 2.0 m, 50 cm, and 200 mm.
\item True or False: CBr 4 has a tetrahedral molecular geometry due to its four equivalent bonding pairs of electrons around the central carbon atom.
\item True or False: The conjugate acid of the H$_{2}$PO$_{4}$- ion is H$_{4}$PO$_{4}$-.
\item True or False: The frequency of light emitted by a helium-neon laser at 632.8 nm is 5.750 x $10^14$ Hz.
\end{enumerate}

\section*{Long Form Response (20 marks)}
\textit{Answer each question with a clear and complete explanation.}
\begin{enumerate}[label=\arabic*.]
\setcounter{enumi}{10}
\item Calculate the mass of calcium hydroxide required to neutralize 28.0 g of hydrochloric acid, and then determine the mass of phosphoric acid that this amount of calcium hydroxide can neutralize. Discuss the underlying chemical reactions and stoichiometric principles involved in your answer.
\item Discuss the relationship between the chemical potentials of components in a binary mixture of water and ethanol, given a mole fraction of water at 0.60 and a change in water's chemical potential. How does this affect the chemical potential of ethanol?
\end{enumerate}

\vfill
\noindent\textit{End of Paper}
\end{document}